% Report-ready component excerpt. Requires amsmath, amssymb and booktabs.
\section{Compact equations from polynomial centers}

This component retains ten strictly positive unknowns and the complete
four-integer coefficient tuple of the inspected native construction.
It composes the explicit trace of package~163 with the inverse Lucas
correction of package~164, resolves the former package's stated open
$b^2$ transport case, and introduces polynomial centers that lower the
degree of the final norm. The following are literal ordinary polynomials;
the intermediate symbols are abbreviations, not additional unknowns.

\begin{center}
\begin{tabular}{lrrrr}
\toprule
Variant & Operations & Exact degree & Weight $L$ & Root scale $\lambda$\\
\midrule
Trace and inverse correction &85&232&$bV$&1\\
Square transport &85&232&$b^2$&1\\
Lower center &87&228&$bV$&1\\
Symmetric center &88&224&$bV$&1\\
Quotient center &96&220&$bV$&2\\
\bottomrule
\end{tabular}
\end{center}

\subsection{Literal formula}

Let $J,F,\alpha,i,T,w,s,n,\eta,\zeta>0$ be the independent integer
variables. The fixed integers are $\beta,K_0,M_C,M_F$ from the whole
directly indexed BFS943/radix67 tuple. Define
\begin{align*}
 h&=\beta J,&q&=h+1,& C&=q-F-\alpha,\\
 X&=qw,&Y&=qs,&k&=\eta+\zeta,&c&=kY+\eta,\\
 A&=Y(X+1)+1,&\Delta&=A^2-1,&H&=4A-5,\\
 \ell&=(XY)(kY),&E&=1+\ell(\ell+k),\\
 R&=[h(q-F)+\alpha](q^2-1)+(M_C+qM_F)J,\\
 B&=(K_0+w)C-F,&U&=ic^2,&a&=\Delta U^2E,\\
 F_0&=1+a,&G&=EF_0,&f&=1+2a,&Q&=4\Delta aF_0,\\
 b&=cT-R,&V&=fb-c,&K&=X(3c-2X)+2,\\
 e_6&=6(k-R-1)-\ell(k^2-4),&D&=6q^6,&M&=hDH.
\end{align*}
For the quotient center also put
\[
 Q_E=(kY)\bigl[\ell+k-(kY)(Y+1)\bigr].
\]
The centers $C_*$ in the table's order are
\[
 0,\qquad0,\qquad(A-1)UE,\qquad AUE,
 \qquad 2AUE-UQ_E.
\]
For the corresponding $L$ and $\lambda$, set
\begin{align*}
 Z&=hD\bigl[H(n-LB-C_*V)+K\bigr]
       -H\bigl[DB+h(e_6-6q^4s)\bigr],\\
 P&=Q\bigl(Z^2-\lambda^2M^2GV^2\bigr)
       +\bigl(\lambda^2M^2G-Z^2\bigr).
\end{align*}
The single equation is $P=0$. Each displayed exponent is a fixed integer.
The literal graphs use only binary addition, subtraction and multiplication.

\subsection{Why the new centers preserve the calibration}

At any positive zero, the inherited root extraction and signed folding
give $G=g^2$ and an integer $y>0$ satisfying
\[
 QV^2-(Q-1)y^2=1,\qquad Z/M=\pm\lambda gy.
\]
The strong trace and the exclusion of the auxiliary indices one and three
give, before the signs of $B$ and $R$ are known,
\[
 g<f/6,\qquad |b|>8f^3-\tfrac12>4f^2,
 \qquad y<\tfrac32|V|,\qquad |V|<\tfrac32f|b|.
\]
The two unscaled nonzero centers lie between zero and $3g/2$; the scaled
center lies between zero and $3g$. Thus every selected variant satisfies
\[
 \lambda gy+|C_*V|<3L/8.
\]
This bound treats both signs of $V$.

For proof notation only, put
\[
 \rho=\frac{k-R-1-XY^2k(k^2-4)/6}{q^6},\quad
 \sigma=\frac{s}{q^2},\quad
 \kappa=\frac KH,\quad t=\frac Bh,\quad W=1+hL.
\]
The inverse correction gives $\rho<|b|/q^6$ uniformly in the integer $R$.
The carrier identity is
\[
 Z/M=n+\kappa+\sigma-Wt-\rho-C_*V.
\]
If $B\le-1$, then $Wt<-L$, so
\[
 n<3L/8-L+|b|/q^6<3L/8-L+L/2<0,
\]
a contradiction. This establishes the positive native branch without
assuming a rank sign or canceling a denominator prematurely.

The center $C_*V$ is an integer polynomial. Consequently the subsequent
index congruences and the projection-before-denominator argument are
unchanged: first recover the main index, then $X=2^R$, then the dyadic
radix, and only then the full divisibilities. Conversely, the retained
canonical witnesses satisfy
$g>\kappa+\sigma+|\rho|$. Choosing
\[
 n=\lambda gy+C_*V+(1+hL)t+\rho-\kappa-\sigma
\]
gives a strictly positive integer and makes the new polynomial vanish.
With the named inherited native/refutation calibration this proves in ZFC
that each polynomial has a positive zero exactly when there is a finite
syntactic ZFC refutation. Negation gives the requested consistency
equivalence; consistency itself is not asserted.

\subsection{Exact cancellation and degree}

Let $Z_0$ be the carrier with its center removed, put
$\overline G=\lambda^2G$, and $S=C_*^2-\overline G$. The free polynomial
identity
\[
 P=(Q-1)Z_0^2-2(Q-1)MC_*VZ_0
   +M^2\bigl[(Q-1)S-\overline G\bigr]V^2+M^2\overline G
\]
exposes the degree cancellation. For the three nonzero centers,
\begin{align*}
 ((A-1)UE)^2-G&=-2(A-1)(UE)^2-E,\\
 (AUE)^2-G&=(UE)^2-E,\\
 (2AUE-UQ_E)^2-4G&=4U^2ER_E+U^2Q_E^2-4E,
\end{align*}
where the exact positive-remainder identity is
\[
 E=AQ_E+R_E,\qquad
 R_E=1+k^2Y(Y+1)\bigl(Y(Y+1)-1\bigr)>0.
\]
The three defects have degrees $46,42,38$, respectively. In the literal
ten-variable polynomial,
\[
 \deg(Q,M,G,V,UE,Z_0)=(80,11,50,40,21,57).
\]
The unique dominant terms in the centered outputs therefore have degrees
$228,224,220$.

For complete leading components, let
$q_0=\beta J$, $k_0=\eta+\zeta$, and
$B_0=k_0sT-q_0^2(q_0-F)$. Then
\begin{align*}
 [P_{232}]_{232}&=-9\cdot2^{10}q_0^{100}w^{26}s^{58}i^{10}k_0^{32}B_0^2,\\
 [P_{228}]_{228}&=-9\cdot2^{11}q_0^{98}w^{25}s^{57}i^{10}k_0^{32}B_0^2,\\
 [P_{224}]_{224}&= 9\cdot2^{10}q_0^{96}w^{24}s^{56}i^{10}k_0^{32}B_0^2,\\
 [P_{220}]_{220}&= 9\cdot2^{12}q_0^{94}w^{22}s^{56}i^{10}k_0^{32}B_0^2.
\end{align*}
These are nonzero for the actual positive $\beta$ and certify exact total
degrees. The full tie between the top terms of $cT$ and $R$ is retained.

\subsection{Constants, costs and finite checks}

The four fixed numeral payloads have the following bit lengths:
\begin{center}
\begin{tabular}{lrrrr}
\toprule
Port &$\beta$&$K_0$&$M_C$&$M_F$\\
\midrule
Bits&590069&592884&590069&590070\\
\bottomrule
\end{tabular}
\end{center}
All raw bytes are supplied and reconstructed from the pinned finite recipe.
The bit lengths measure the supplied numeral ports, not the height of the
fully expanded final polynomial. Arithmetic counts treat these fixed
integers as constants and count every live binary operation thereafter.

Sharing $V$ through
$n-bVB-C_*V=n-V(bB+C_*)$ saves one multiplication. Hence the three
centered costs are exactly $87$, $88$ and $96$. The independent checker
covers all five graphs with 199 free-polynomial cuts, authenticates the
four actual numeral payloads and 35 inherited source files, verifies the
complete leading components, and rejects 73 deliberate operator
corruptions. Normal and optimized Python receipts agree. The separate
finite recipe reconstructs all four integers, including the required native
$M_F$ shift. These are artifact and algebra checks; the unbounded arithmetic
proof and inherited logical calibration remain explicit dependencies.
